\section{Conclusion}\label{sec:conclusion}

The objective was a protocol under which discourse becomes something of value to those who took part in it and to
their community. It needs a unit of knowledge whose use is a fact, and the construction is such a unit: where
applying it to a text is deterministic, a use can be re-executed by a stranger, and a network can pay without
paying for anyone's opinion. It needs also a measure of worth that a few parties cannot manufacture, and the
measure proposed is breadth: how many independent communities have taken a construction up and kept it.

The protocol follows from those two. Communities publish their grammars for open use and signal what they are for.
A community that cannot read something looks it up, tries what it finds, and adopts what goes on working.
Adoption is free and is never pushed. An originator is paid according to the breadth of its constructions, with
nothing below a threshold, and a share passes to the parents each construction declares. Membership is free, and
members whose contributions become constructions earn stake in what they made. Nodes supply computation, and the
stake they pledge sets a ceiling that only breadth can fill, so capital follows uptake.

What cannot be computed is whether the communities that hold a construction are separate parties. That is anchored
in stake. Under the assumptions of Section~\ref{sec:security}, the ceiling rate is set no higher than the cost of
capital plus the chance of being caught posing as many, and within that limit a community that wants to earn has
one course open: to make constructions that others, who owe it nothing, find they need and keep.
